\phantomsection
\addcontentsline{lot}{section}{\protect\numberline{}Tabel 49 - Status gizi balita berdasarkan indeks BB/U, TB/U, dan BB/TB}
\ra{1.3}
%Table generated by Excel2LaTeX from sheet '48'

\begin{tabular}{rY{9em}Y{8em}rZ{3em}rrrrrrrrr}
    \multicolumn{14}{l}{Tabel 49}\\
    \multicolumn{14}{c}{STATUS GIZI BALITA BERDASARKAN INDEKS BB/U, TB/U, DAN BB/TB}\\
	\multicolumn{14}{c}{DAN TATA LAKSANA MENURUT KECAMATAN DAN PUSKESMAS}\\
    \multicolumn{14}{c}{KABUPATEN BELITUNG TIMUR}\\
    \multicolumn{14}{c}{TAHUN \tP}\\
    \toprule
    \multirow{2}[0]{*}{NO} & \multirow{2}[0]{*}{KECAMATAN} & \multirow{2}[0]{*}{PUSKESMAS} & \multirow{2}[0]{5em}{\raggedleft JUMLAH BALITA YANG DITIMBANG} & \multicolumn{2}{X{7em}}{BALITA BERAT BADAN KURANG (BB/U)} & \multirow{2}[0]{7em}{\raggedleft JUMLAH BALITA YANG DIUKUR TINGGI BADAN} & \multicolumn{2}{X{6em}}{BALITA PENDEK (TB/U)} & \multirow{2}[0]{5em}{\raggedleft JUMLAH BALITA YANG DIUKUR} & \multicolumn{2}{X{6em}}{BALITA GIZI KURANG (BB/TB : $<$ -2 s.d -3 SD)} & \multicolumn{2}{X{6em}}{BALITA GIZI BURUK (BB/TB: $<$ -3 SD)}\\
%    & & & & & & & & & & & & & \\ % extra row to accomodate some multirow columns above
    \cmidrule(l{2pt}r{2pt}){5-6}\cmidrule(l{2pt}r{2pt}){8-9}\cmidrule(l{2pt}r{2pt}){11-12}\cmidrule(l{2pt}r{2pt}){13-14}
    & & & & Jml & \% & & Jml & \% & & Jml & \% & Jml & \%\\
    \midrule
    \emph{1} & \emph{2} & \emph{3} & \emph{4} & \emph{5} & \emph{6} & \emph{7} & \emph{8} & \emph{9} & \emph{10} & \emph{11} & \emph{12} & \emph{13} & \emph{14}\\
    \midrule
	1 & Manggar           & Manggar       &       &     &       &       &     &       &       &     &      &   &      \\
	2 & Damar             & Mengkubang    &       &     &       &       &     &       &       &     &      &   &      \\
	3 & Kelapa Kampit     & Kelapa Kampit &       &     &       &       &     &       &       &     &      &   &      \\
	4 & Gantung           & Gantung       &       &     &       &       &     &       &       &     &      &   &      \\
	5 & Simpang Renggiang & Renggiang     &       &     &       &       &     &       &       &     &      &   &      \\
	6 & Simpang Pesak     & Simpang Pesak &       &     &       &       &     &       &       &     &      &   &      \\
	7 & Dendang           & Dendang       &       &     &       &       &     &       &       &     &      &   &      \\
    \midrule
    \multicolumn{3}{l}{JUMLAH}            & 7.524 & 618 &  8,21 & 7.523 & 347 &  4,61 & 7.523 & 239 & 3,18 & 1 & 0,01 \\
    \bottomrule
\end{tabular}%


\begin{tabular}{rY{9em}Y{8em}rZ{5em}rrZ{5em}r}
    \multicolumn{9}{l}{Tabel 49 (lanj.}\\
    \toprule
    \multirow{2}[0]{*}{NO} & \multirow{2}[0]{*}{KECAMATAN} & \multirow{2}[0]{*}{PUSKESMAS} & \multirow{3}[0]{9em}[-1ex]{\raggedleft JUMLAH BALITA USIA 6-59 BULAN YANG GIZI KURANG DENGAN ATAU TANPA STUNTING} & \multicolumn{2}{X{10em}}{BALITA GIZI KURANG MENDAPAT MAKANAN TAMBAHAN} & \multirow{2}[0]{7em}[-1ex]{\raggedleft JUMLAH KASUS BALITA GIZI BURUK} & \multicolumn{2}{X{9em}}{BALITA GIZI BURUK MENDAPAT TATALAKSANA} \\
    \cmidrule(l{2pt}r{2pt}){5-6}\cmidrule(l{2pt}r{2pt}){8-9}
    & & & & Jml & \% & & Jml & \% \\
    \midrule
    \emph{1} & \emph{2} & \emph{3} & \emph{15} & \emph{16} & \emph{17} & \emph{18} & \emph{19} & \emph{20} \\
    \midrule
	1 & Manggar           & Manggar       &       &     &       &       &     &       \\
	2 & Damar             & Mengkubang    &       &     &       &       &     &       \\
	3 & Kelapa Kampit     & Kelapa Kampit &       &     &       &       &     &       \\
	4 & Gantung           & Gantung       &       &     &       &       &     &       \\
	5 & Simpang Renggiang & Renggiang     &       &     &       &       &     &       \\
	6 & Simpang Pesak     & Simpang Pesak &       &     &       &       &     &       \\
	7 & Dendang           & Dendang       &       &     &       &       &     &       \\
    \midrule
    \multicolumn{3}{l}{JUMLAH}            & 7.524 & 618 &  8,21 & 7.523 & 347 &  4,61 \\
    \bottomrule
\end{tabular}%
    
\vfill
Sumber: Subkoordinator Kesehatan Keluarga dan Gizi\par 
